\chapter{The Rauch Comparison Theorem}
\label{chap:dc-rauch-comparison-theorem}

Faithful transcription of do Carmo, \emph{Riemannian Geometry}. Each numbered
statement keeps do Carmo's number, kind and (name); proofs keep his explicit
cross-references ("by \cref{lem:dc-ch10-2-2}", "\cref{prop:dc-ch5-3-6}", "Chapter 9")
so the dependency graph is recoverable from the text.

\section{Introduction}

Let $M$ be a Riemannian manifold. As we saw in Chapter 5, if
$\gamma:[0,\ell]\to M$ is a normalized geodesic and $J$ is a Jacobi field along
$\gamma$ with $J(0)=0$, $|J'(0)|=1$ and $\langle J'(0),\gamma'(0)\rangle=0$, then

$$
|J(t)|=t-\frac{K}{6}t^3+R,\qquad\lim_{t\to 0}\frac{R}{t^3}=0,
$$

where $K$ is the sectional curvature at $\gamma(0)$ with respect to the plane
generated by $\gamma'(0)$ and $J'(0)$. Therefore, if $t$ is small, the smaller
$K$ is, the larger $|J(t)|$ will be. Consider now another Riemannian manifold
$\widetilde{M}$, a geodesic $\tilde\gamma:[0,\ell]\to\widetilde{M}$ and a Jacobi
field $\tilde J$ along $\tilde\gamma$ satisfying $\tilde J(0)=0$, $|\tilde J'(0)|=1$,
$\langle\tilde J'(0),\tilde\gamma'(0)\rangle=0$. Suppose that

$$
\widetilde{K}(\tilde\gamma'(0),\tilde J'(0))\geq K(\gamma'(0),J'(0)).
$$

It follows that, for small $t$, $|\tilde J(t)|\leq|J(t)|$.

The Theorem of Rauch furnishes conditions so that the inequality above is valid
without the restriction that $t$ be small. More precisely, the theorem asserts
(see the statement in Section 2) that if $\tilde\gamma:[0,\ell]\to\widetilde{M}$
does not have conjugate points and

$$
\widetilde{K}(\tilde\gamma'(t),\tilde J(t))\geq K(\gamma'(t),J(t)),\quad t\in(0,\ell),
$$

then $|\tilde J(t)|\leq|J(t)|$. Intuitively, Rauch's Theorem expresses the
plausible fact that as the curvature grows, lengths shorten.

In dimension two, such a theorem is an easy consequence of the classical theorem
of Sturm on ordinary differential equations. In fact, the Jacobi equations for
$J=fe_2$, $\tilde J=\tilde f\tilde e_2$ ($e_2(t)$ and $\tilde e_2(t)$ unit
parallel vector fields normal to $\gamma(t)$ and $\tilde\gamma(t)$) can be
written

$$
f''(t)+K(t)f(t)=0,\quad f(0)=0,\quad t\in[0,\ell],
$$

$$
\tilde f''(t)+\widetilde{K}(t)\tilde f(t)=0,\quad\tilde f(0)=0,\quad t\in[0,\ell].
$$

Sturm's Theorem asserts that if $f'(0)=\tilde f'(0)>0$, $\tilde f(t)\neq 0$ on
$(0,\ell]$, and $\widetilde{K}(t)\geq K(t)$, then $\tilde f(t)\leq f(t)$.

In dimension higher than two the proof is much less simple; a presentation of
the theorem was made for the first time in 1951 by Rauch. The proof presented in
Section 2 is an elaboration due to various mathematicians in the fifties, among
whom we should mention Ambrose and Singer. In Section 3 we make an application of
the ideas in the Rauch Theorem to the theory of isometric immersions. In Section
4 we introduce the notion of focal point, which generalizes the idea of
conjugate point, and we extend the Theorem of Rauch to this situation.

\section{The Theorem of Rauch}

In the proof of Rauch's Theorem, we need a certain number of facts that we
establish in the form of lemmas. The Index Lemma (\cref{lem:dc-ch10-2-2}) is a basic fact
that has application to many other situations. The following lemma is a
particular case of \cref{lem:dc-ch0-5-5}.

\begin{lemma}[Lemma 2.1]
  \label{lem:dc-ch10-2-1}
  \dcref{ch10:2.1}
  Let $h:[0,1]\to\mathbb{R}$ be a differentiable function with $h(0)=0$. Then there
  exists a differentiable function $\phi:[0,1]\to\mathbb{R}$, with
  $\phi(0)=\frac{dh}{dt}(0)$, $h(t)=t\phi(t)$, $t\in[0,1]$.
  
  Let $M$ be a Riemannian manifold and let $\gamma:[0,a]\to M$ be a geodesic of
  $M$. Let $V$ be a piecewise differentiable vector field along $\gamma$. For all
  $t_o\in[0,a]$, we write
  
  $$
  \int_0^{t_o}\{\langle V',V'\rangle-\langle R(\gamma',V)\gamma',V\rangle\}\,dt=I_{t_o}(V,V).
  $$
\end{lemma}

\begin{lemma}[The Index Lemma]
  \label{lem:dc-ch10-2-2}
  \dcref{ch10:2.2}
  \uses{lem:dc-ch10-2-1}
  Let $\gamma:[0,a]\to M$ be a geodesic without conjugate points to $\gamma(0)$ in
  the interval $(0,a]$. Let $J$ be a Jacobi field along $\gamma$, with
  $\langle J,\gamma'\rangle=0$, and let $V$ be a piecewise differentiable vector
  field along $\gamma$, with $\langle V,\gamma'\rangle=0$. Suppose that
  $J(0)=V(0)=0$ and that $J(t_o)=V(t_o)$, $t_o\in(0,a]$. Then
  
  $$
  I_{t_o}(J,J)\leq I_{t_o}(V,V),
  $$
  
  and equality occurs if and only if $V=J$ on $[0,t_o]$.
\end{lemma}
\begin{proof}
  The vector space $\mathcal{J}$ of Jacobi fields $J$ along $\gamma$ with
  $J(0)=0$ and $\langle J,\gamma'\rangle=0$ has dimension $n-1$, where $n=\dim M$.
  Let $\{J_1,\dots,J_{n-1}\}$ be a basis for this space; then $J=\sum_i\alpha_i J_i$
  with constants $\alpha_i$. Since there are no conjugate points in $(0,a]$, for
  all $t\neq 0$ the vectors $J_1(t),\dots,J_{n-1}(t)$ form a basis of the
  orthogonal complement of $\gamma'(t)$, so for $t\neq 0$ we can write
  $V(t)=\sum_i f_i(t)J_i(t)$, with $f_i$ piecewise differentiable on $(0,a]$. Using
  \cref{lem:dc-ch10-2-1} (write $J_i(t)=tA_i(t)$, so the $A_i(0)=J_i'(0)$ are linearly
  independent, and $V(t)=\sum_i g_i(t)A_i(t)$ with $g_i(0)=0$, then $g_i(t)=th_i(t)$
  with $f_i=h_i$), the $f_i$ extend piecewise differentiably to $[0,a]$.
  
  One shows that, on the interior of each subinterval where $f_i$ is
  differentiable,
  
  $$
  \langle V',V'\rangle-\langle R(\gamma',V)\gamma',V\rangle=\Big\langle\sum_i f_i'J_i,\sum_j f_j'J_j\Big\rangle+\frac{d}{dt}\Big\langle\sum_i f_iJ_i,\sum_j f_jJ_j'\Big\rangle.\qquad(1)
  $$
  
  Here one uses $R(\gamma',V)\gamma'=-\sum_i f_iJ_i''$, and the identity
  
  $$
  \Big\langle\sum_i f_iJ_i',\sum_j f_j'J_j\Big\rangle=\Big\langle\sum_i f_iJ_i,\sum_j f_j'J_j'\Big\rangle,\qquad(2)
  $$
  
  which follows from $h(t)=\langle J_i',J_j\rangle-\langle J_i,J_j'\rangle$ having
  $h(0)=0$ and $h'(t)=-\langle R(\gamma',J_i)\gamma',J_j\rangle+\langle J_i,R(\gamma',J_j)\gamma'\rangle=0$,
  so $h\equiv 0$. Applying (1) to $V$ and $J$,
  
  $$
  I_{t_o}(V,V)=\Big\langle\sum_i f_iJ_i,\sum_j f_jJ_j'\Big\rangle(t_o)+\int_0^{t_o}\Big\langle\sum_i f_i'J_i,\sum_j f_j'J_j\Big\rangle dt,
  $$
  
  $$
  I_{t_o}(J,J)=\Big\langle\sum_i\alpha_iJ_i,\sum_j\alpha_jJ_j'\Big\rangle(t_o).
  $$
  
  Because $J(t_o)=V(t_o)$, we have $\alpha_i=f_i(t_o)$, hence
  
  $$
  I_{t_o}(V,V)=I_{t_o}(J,J)+\int_0^{t_o}\Big|\sum_i f_i'J_i\Big|^2 dt.\qquad(3)
  $$
  
  It follows from (3) that $I_{t_o}(V,V)\geq I_{t_o}(J,J)$. If equality holds, then
  $\sum_i f_i'J_i=0$; since the $J_i$ are linearly independent for $t\neq 0$, by
  continuity $f_i'=0$ for all $i$ on $[0,t_o]$. Therefore $f_i$ is constant, and
  since $f_i(t_o)=\alpha_i$, we have $f_i(t)=\alpha_i$, that is, $V=J$. $\blacksquare$
  
  We are now in a position to prove Rauch's Theorem. In what follows $M^n$ denotes
  a manifold of dimension $n$.
\end{proof}

\begin{theorem}[Rauch]
  \label{thm:dc-ch10-2-3}
  \dcref{ch10:2.3}
  \uses{lem:dc-ch10-2-2, prop:dc-ch5-3-6}
  Let $\gamma:[0,a]\to M^n$ and $\tilde\gamma:[0,a]\to\widetilde{M}^{n+k}$, $k\geq 0$,
  be geodesics with the same velocity (i.e. $|\gamma'(t)|=|\tilde\gamma'(t)|$), and
  let $J$ and $\tilde J$ be Jacobi fields along $\gamma$ and $\tilde\gamma$,
  respectively, such that
  
  $$
  J(0)=\tilde J(0)=0,\quad\langle J'(0),\gamma'(0)\rangle=\langle\tilde J'(0),\tilde\gamma'(0)\rangle,\quad|J'(0)|=|\tilde J'(0)|.
  $$
  
  Assume that $\tilde\gamma$ does not have conjugate points on $(0,a]$ and that, for
  all $t$ and all $x\in T_{\gamma(t)}(M)$, $\tilde x\in T_{\tilde\gamma(t)}(\widetilde{M})$,
  we have
  
  $$
  \widetilde{K}(\tilde x,\tilde\gamma'(t))\geq K(x,\gamma'(t)),
  $$
  
  where $K(x,y)$ denotes the sectional curvature with respect to the plane
  generated by $x$ and $y$. Then
  
  $$
  |\tilde J|\leq|J|.
  $$
  
  In addition, if for some $t_o\in(0,a]$ we have $|\tilde J(t_o)|=|J(t_o)|$, then
  $\widetilde{K}(\tilde J(t),\tilde\gamma'(t))=K(J(t),\gamma'(t))$ for all $t\in[0,t_o]$.
\end{theorem}
\begin{proof}
  By \cref{prop:dc-ch5-3-6}, the condition
  $\langle J'(0),\gamma'(0)\rangle=\langle\tilde J'(0),\tilde\gamma'(0)\rangle$ is
  equivalent (with $J(0)=\tilde J(0)=0$) to $\langle J,\gamma'\rangle=\langle\tilde J,\tilde\gamma'\rangle$;
  since $\langle J,\gamma'\rangle\gamma'=\langle J'(0),\gamma'(0)\rangle t\gamma'+\langle J(0),\gamma'(0)\rangle\gamma'$,
  the tangential components of $J$ and $\tilde J$ have the same length, so we may
  suppose $\langle J,\gamma'\rangle=0=\langle\tilde J,\tilde\gamma'\rangle$.
  
  If $|J'(0)|=|\tilde J'(0)|=0$, then $|J|=|\tilde J|=0$. Otherwise put
  $|J(t)|^2=v(t)$ and $|\tilde J(t)|^2=\tilde v(t)$. Since $\tilde\gamma$ has no
  conjugate points on $(0,a]$, $v/\tilde v$ is well-defined on $(0,a]$, and by
  L'Hospital's rule $\lim_{t\to 0}v(t)/\tilde v(t)=1$. To prove $|\tilde J|\leq|J|$
  it is enough to prove $\frac{d}{dt}(v/\tilde v)\geq 0$, i.e. $v'\tilde v\geq v\tilde v'$.
  Fixing $t_o\in(0,a]$ with $v(t_o)\neq 0$, put $U(t)=\frac{1}{\sqrt{v(t_o)}}J(t)$,
  $\tilde U(t)=\frac{1}{\sqrt{\tilde v(t_o)}}\tilde J(t)$; one computes
  
  $$
  \frac{v'(t_o)}{v(t_o)}=\langle U,U\rangle'(t_o)=2\int_0^{t_o}\{\langle U',U'\rangle-\langle U,R(\gamma',U)\gamma'\rangle\}dt=2I_{t_o}(U,U),
  $$
  
  and similarly $\tilde v'(t_o)/\tilde v(t_o)=2I_{t_o}(\tilde U,\tilde U)$. By the
  arbitrariness of $t_o$, it suffices to prove $I_{t_o}(\tilde U,\tilde U)\leq I_{t_o}(U,U)$.
  
  Let $\{e_1,\dots,e_n\}$ and $\{\tilde e_1,\dots,\tilde e_{n+k}\}$ be parallel
  orthonormal bases along $\gamma$ and $\tilde\gamma$ with $e_1(t)=\gamma'(t)/|\gamma'(t)|$,
  $e_2(t_o)=U(t_o)$, $\tilde e_1(t)=\tilde\gamma'(t)/|\tilde\gamma'(t)|$,
  $\tilde e_2(t_o)=\tilde U(t_o)$. To each $V(t)=\sum_i g_i(t)e_i(t)$ along $\gamma$
  associate $\phi V=\sum_{i=1}^n g_i(t)\tilde e_i(t)$ along $\tilde\gamma$. Then
  
  $$
  \langle\phi V_1,\phi V_2\rangle=\langle V_1,V_2\rangle,\qquad(4)
  $$
  
  $$
  (\phi V)'=\phi(V').\qquad(5)
  $$
  
  From the hypothesis on the curvature and equal velocity, $I_{t_o}(\phi(U),\phi(U))\leq I_{t_o}(U,U)$.
  Since $\tilde U$ and $\phi(U)$ satisfy the hypothesis of \cref{lem:dc-ch10-2-2} and $\tilde U$
  is a Jacobi field, \cref{lem:dc-ch10-2-2} gives
  
  $$
  I_{t_o}(\tilde U,\tilde U)\leq I_{t_o}(\phi(U),\phi(U))\leq I_{t_o}(U,U),
  $$
  
  which proves the inequality.
  
  If $|J(t_o)|=|\tilde J(t_o)|$ for some $t_o$, then $v'\tilde v=v\tilde v'$ on
  $(0,t_o]$, so the inequalities above are equalities, i.e.
  $I_t(\phi(U),\phi(U))=I_t(U,U)$ on $(0,t_o]$. Since $\phi$ satisfies (4) and (5),
  we conclude from the equality and the curvature hypothesis that
  $K(J(t),\gamma'(t))=\widetilde{K}(\tilde J(t),\tilde\gamma'(t))$ for $t\in(0,t_o]$,
  hence by continuity on $[0,t_o]$. $\blacksquare$
  
  An immediate application of Rauch's Theorem allows us to obtain information on the
  location of conjugate points from bounds on the curvature.
\end{proof}

\begin{proposition}[Proposition 2.4]
  \label{prop:dc-ch10-2-4}
  \dcref{ch10:2.4}
  Suppose that the sectional curvature $K$ of a Riemannian manifold $M$ satisfies
  
  $$
  0<L\leq K\leq H,
  $$
  
  where $H$ and $L$ are constants. Let $\gamma$ be a geodesic in $M$. Then the
  distance $d$ along $\gamma$ between two consecutive conjugate points of $\gamma$
  satisfies
  
  $$
  \frac{\pi}{\sqrt{H}}\leq d\leq\frac{\pi}{\sqrt{L}}.
  $$
\end{proposition}
\begin{proof}
  To get $d\geq\pi/\sqrt{H}$, compare $M^n$ with the sphere $S^n(H)$ of
  curvature $H$. Let $\gamma:[0,\ell]\to M$ be a normalized geodesic with
  $\gamma(0)=p$ and $J$ a Jacobi field with $J(0)=0$, $\langle J,\gamma'\rangle=0$.
  Choose $\tilde p\in S^n(H)$ and a normalized geodesic $\tilde\gamma$ with
  $\tilde\gamma(0)=\tilde p$, and $\tilde J$ with $\tilde J(0)=0$,
  $\langle\tilde J,\tilde\gamma'\rangle=0$, $|\tilde J'(0)|=|J'(0)|$. Since
  $\tilde\gamma$ has no conjugate points on $(0,\pi/\sqrt{H})$, by Rauch's Theorem
  $|J(t)|\geq|\tilde J(t)|>0$ on $(0,\pi/\sqrt{H})$. Therefore $d\geq\pi/\sqrt{H}$.
  To obtain $d\leq\pi/\sqrt{L}$, make the analogous comparison with $S^n(L)$: if
  $d>\pi/\sqrt{L}$, Rauch's Theorem applies and $S^n(L)$ would have all its
  conjugate points after $\pi/\sqrt{L}$, which is absurd. $\blacksquare$
  
  One of the typical applications of Rauch's Theorem consists in estimating the
  lengths of curves in a Riemannian manifold in which we can estimate the
  curvature.
\end{proof}

\begin{proposition}[Proposition 2.5]
  \label{prop:dc-ch10-2-5}
  \dcref{ch10:2.5}
  Let $M^n$ and $\widetilde{M}^n$ be Riemannian manifolds and suppose that for all
  $p\in M$, $\tilde p\in\widetilde{M}$, $\sigma\subset T_pM$, $\tilde\sigma\subset T_{\tilde p}\widetilde{M}$,
  we have $\widetilde{K}_{\tilde p}(\tilde\sigma)\geq K_p(\sigma)$. Let $p\in M$,
  $\tilde p\in\widetilde{M}$ and fix a linear isometry $i:T_pM\to T_{\tilde p}\widetilde{M}$.
  Let $r>0$ be such that the restriction $\exp_p|_{B_r(0)}$ is a diffeomorphism and
  $\exp_{\tilde p}|_{\tilde B_r(0)}$ is non-singular. Let $c:[0,a]\to\exp_p(B_r(0))\subset M$
  be a differentiable curve and define $\tilde c:[0,a]\to\exp_{\tilde p}(\tilde B_r(0))\subset\widetilde{M}$
  by
  
  $$
  \tilde c(s)=\exp_{\tilde p}\circ i\circ\exp_p^{-1}(c(s)),\quad s\in[0,a].
  $$
  
  Then $\ell(c)\geq\ell(\tilde c)$.
\end{proposition}
\begin{proof}
  Consider $\bar c(s)=\exp_p^{-1}c(s)$ in $T_pM$. For $s$ fixed, the radial
  geodesic $\gamma_s(t)=\exp_p t\bar c(s)$ gives the parametrized surface
  $f(t,s)=\gamma_s(t)$, $0\leq s\leq a$, $0\leq t\leq 1$. Then $J_s(t)=\frac{\partial f}{\partial s}(t,s)$
  is a Jacobi field along $\gamma_s$ with $J_s(0)=0$, $J_s(1)=c'(s)$, and
  $\frac{DJ_s}{dt}(0)=\bar c'(s)$. Consider in $\widetilde{M}$ the surface
  $\tilde f(t,s)=\exp_{\tilde p}ti(\bar c(s))=\tilde\gamma_s(t)$; then $\tilde J_s(t)=\frac{\partial\tilde f}{\partial s}(t,s)$
  is a Jacobi field with $\tilde J_s(0)=0$, $\tilde J_s(1)=\tilde c'(s)$,
  $\frac{D\tilde J_s}{dt}(0)=i\bar c'(s)$. Because $i$ is an isometry,
  $|J_s(0)|=|\tilde J_s(0)|=0$, $|J_s'(0)|=|\tilde J_s'(0)|$ and
  $\langle\tilde J_s'(0),\tilde\gamma_s'(0)\rangle=\langle\bar c'(s),\gamma_s'(0)\rangle=\langle J_s'(0),\gamma_s'(0)\rangle$.
  Applying Rauch's Theorem,
  
  $$
  |\tilde c'(s)|=|\tilde J_s(1)|\leq|J_s(1)|=|c'(s)|,
  $$
  
  hence, by integration, $\ell(c)\geq\ell(\tilde c)$. $\blacksquare$
\end{proof}

\begin{remark}[Remark 2.6]
  \label{rem:dc-ch10-2-6}
  \dcref{ch10:2.6}
  The methods used in Rauch's Theorem can be employed to obtain an upper bound for
  the volume of a normal ball in a Riemannian manifold whose Ricci curvature is
  bounded below. More precisely, let $M$ be a complete Riemannian manifold with
  $\mathrm{Ric}_M\geq H$, and let $B_r(p)\subset M$, $B_r(\tilde p)\subset\widetilde{M}(H)$
  be normal balls of radius $r$, where $\widetilde{M}(H)$ is the complete, simply
  connected manifold of constant sectional curvature $H$. Then
  
  $$
  \mathrm{vol}(B_r(p))\leq\mathrm{vol}\,B_r(\tilde p),
  $$
  
  and, if equality holds, $\mathrm{Ric}_M(\gamma'(t))\equiv H$ for all $t<r$ and all
  normalized radial geodesics $\gamma(t)$ in $B_r(p)$ (see Bishop and Crittenden,
  \emph{Geometry of Manifolds}, Academic Press, 1964, p. 256, and K. Tenenblat, *On the
  Rauch comparison theorem for volumes*, Bol. Soc. Bras. Mat. 4 (1973), 31–39). In
  contrast to Rauch's Theorem, the theorem above does not extend to the case
  $\mathrm{Ric}_M\leq L$; a counterexample is described in the article cited above
  by K. Tenenblat.
\end{remark}

\begin{remark}[Remark 2.7]
  \label{rem:dc-ch10-2-7}
  \dcref{ch10:2.7}
  \uses{rem:dc-ch9-3-6, rem:dc-ch10-2-6}
  The relation between the Ricci curvature and the volume can be extended to the
  following global theorem: with the same notation as in \cref{rem:dc-ch10-2-6}, suppose that
  $\mathrm{Ric}_M\geq H$ and denote the corresponding volumes by $V_r=\mathrm{vol}\,B_r(p)$
  and $\widetilde{V}_r=\mathrm{vol}\,B_r(\tilde p)$, where $r$ is now an arbitrary
  positive real number. Then for all $R\geq r>0$ we have
  
  $$
  V_r/\widetilde{V}_r\geq V_R/\widetilde{V}_R.
  $$
  
  In addition, the equality occurs for $R\leq\mathrm{diam}(M)$ if and only if
  $B_R(p)\subset M$ is isometric to the ball $B_R(\tilde p)\subset\widetilde{M}(H)$.
  For a proof, see J. Eschenburg, \emph{Comparison theorems and hypersurfaces},
  Manuscripta Math. 59 (1987), 295–323. This reference is used to establish the
  uniqueness theorem mentioned in \cref{rem:dc-ch9-3-6}.
\end{remark}

\section{Applications of the Index Lemma to immersions}

In this Section, we present an application of the Index Lemma to the theory of
isometric immersions. We prove the following theorem, which generalizes previous
theorems of Tompkins, O'Neill, and Chern-Kuiper (see Corollaries 3.4, 3.5 and
Remarks 3.6 and 3.7).

\begin{theorem}[J.D. Moore]
  \label{thm:dc-ch10-3-1}
  \dcref{ch10:3.1}
  \uses{lem:dc-ch10-3-2, lem:dc-ch10-3-3}
  Let $\overline{M}$ be a complete simply connected Riemannian manifold, whose
  sectional curvature satisfies
  
  $$
  \overline{K}\leq b\leq 0.
  $$
  
  Let $M$ be a compact Riemannian manifold whose sectional curvature satisfies
  $K-\overline{K}\leq -b$. If $\dim\overline{M}<2\dim M$, there does not exist an
  isometric immersion $f:M\to\overline{M}$.
\end{theorem}
\begin{proof}
  Suppose such an immersion exists, and obtain a contradiction. Choose
  $\bar p\in\overline{M}$, $\bar p\notin f(M)$, and let $q\in M$ satisfy
  $d(f(q),\bar p)\geq d(f(p),\bar p)$ for all $p\in M$. Let $U\subset M$ be a
  neighborhood of $q$ in which $f$ is an embedding, and identify $U$ with $f(U)$.
  Let $\gamma:[0,\ell]\to\overline{M}$ be the normalized minimizing geodesic with
  $\gamma(0)=\bar p$, $\gamma(\ell)=q$. Using the formula for the first variation,
  $\gamma$ is perpendicular to $U$ at $q$. For all $v\in T_qM$, $|v|=1$, take a
  curve $c(s)$ in $U$ with $c(0)=q$, $c'(0)=v$, and let $\tilde c(s)=\exp_{\bar p}^{-1}(c(s))$.
  The parametrized surface $f(t,s)=\exp_{\bar p}\frac{t}{\ell}\tilde c(s)$ generates
  a Jacobi field $V(t)=\frac{\partial f}{\partial s}(t,0)$ with $V(0)=0$, $V(\ell)=v$;
  the radial geodesics $\gamma_s:t\mapsto f(t,s)$ are shorter than $\gamma$.
  
  We use \cref{lem:dc-ch10-3-2}. Consider a complete manifold $\widetilde{M}(b)$ of constant
  curvature $b$ and dimension $n=\dim\overline{M}$; choosing parallel orthonormal
  bases and the map $\phi$ as before, with $\overline{K}\leq\widetilde{K}\equiv b$,
  we get $I_\ell(V,V)\geq I_\ell(\tilde V,\tilde V)$. Estimating $I_\ell(\tilde V,\tilde V)$
  via the Jacobi field $\tilde J$ along $\tilde\gamma$ with $\tilde J(0)=0$,
  $\tilde J(\ell)=\tilde v$ (Cf. Exercise 4, Chap. 5),
  
  $$
  \tilde J(t)=\frac{\sinh(t\sqrt{-b})}{\sinh(\ell\sqrt{-b})}\tilde w(t)\ (b<0),\qquad\tilde J(t)=\frac{t}{\ell}\tilde w(t)\ (b=0),
  $$
  
  where $\tilde w$ is the parallel vector field along $\tilde\gamma$ with
  $\tilde w(\ell)=\tilde v$. Thus
  one obtains from the proof of the Index Lemma
  
  $$
  I_\ell(\tilde J,\tilde J)=(\coth\ell\sqrt{-b})\sqrt{-b}>\sqrt{-b},
  $$
  
  and, if $b=0$, $I_\ell(\tilde J,\tilde J)>0$. Therefore from the Index Lemma,
  
  $$
  I_\ell(V,V)\geq I_\ell(\tilde V,\tilde V)\geq I_\ell(\tilde J,\tilde J)>\sqrt{-b}.
  $$
  
  Because the geodesics $\gamma_s$ are shorter than $\gamma$,
  $\ell E(\gamma)=(L(\gamma))^2\geq(L(\gamma_s))^2=\ell E(\gamma_s)$, hence
  
  $$
  0\geq\tfrac{1}{2}E''(0)=I_\ell(V,V)+\langle S_{\gamma'(\ell)}V(\ell),V(\ell)\rangle>\sqrt{-b}+\langle S_{\gamma'(\ell)}V(\ell),V(\ell)\rangle.
  $$
  
  Since $V(\ell)=v$, $\langle B(v,v),\gamma'(\ell)\rangle=\langle S_{\gamma'(\ell)}v,v\rangle<-\sqrt{-b}$,
  so for all unit $v\in T_qM$,
  
  $$
  \|B(v,v)\|>\sqrt{-b}.\qquad(6)
  $$
  
  On the other hand, from the Gauss formula, for orthonormal $v,w\in T_qM$,
  
  $$
  \langle B(v,v),B(w,w)\rangle-\|B(v,w)\|^2=K(v,w)-\overline{K}(v,w)\leq -b.\qquad(7)
  $$
  
  That (6) and (7) are incompatible with $\dim\overline{M}<2\dim M$ follows from
  \cref{lem:dc-ch10-3-3}, completing the proof. $\blacksquare$
\end{proof}

\begin{lemma}[Lemma 3.2]
  \label{lem:dc-ch10-3-2}
  \dcref{ch10:3.2}
  Let $E(s)$ be the energy of the curve $t\mapsto f(t,s)$. Then
  
  $$
  \tfrac{1}{2}E''(0)=I_\ell(V,V)+\langle S_{\gamma'(\ell)}V(\ell),V(\ell)\rangle,
  $$
  
  where $S_{\gamma'(\ell)}$ is the linear operator associated to the second
  fundamental form of the immersion $f$ at the point $q$ with respect to the normal
  $\gamma'(\ell)$.
\end{lemma}
\begin{proof}
  We know that (see Remark 2.10, Chap. 9)
  
  $$
  \tfrac{1}{2}E''(0)=I_\ell(V,V)+\Big\langle\frac{D}{ds}\frac{\partial f}{\partial s},\frac{\partial f}{\partial t}\Big\rangle(\ell,0)=I_\ell(V,V)+\langle(\overline{\nabla}_{\overline V}\overline V)(q),\overline N(q)\rangle,
  $$
  
  where $\overline\nabla$ is the covariant derivative of $\overline{M}$, and
  $\overline V$, $\overline N$ are local extensions of $V(\ell)$ and $\gamma'(\ell)$
  chosen so that $\langle\overline V,\overline N\rangle=0$. Therefore
  $\langle\overline\nabla_{\overline V}\overline V,\overline N\rangle(q)=-\langle\overline\nabla_{\overline V}\overline N,\overline V\rangle(q)=\langle V(\ell),S_{\gamma'(\ell)}V(\ell)\rangle$,
  hence the assertion. $\blacksquare$
\end{proof}

\begin{lemma}[Otsuki]
  \label{lem:dc-ch10-3-3}
  \dcref{ch10:3.3}
  Let $B:\mathbb{R}^n\times\mathbb{R}^n\to\mathbb{R}^k$ be a symmetric bilinear form
  such that, for some $b\leq 0$,
  
  $$
  \langle B(v,v),B(w,w)\rangle-\|B(v,w)\|^2\leq -b,\qquad\|B(v,v)\|>\sqrt{-b},
  $$
  
  for all orthonormal pairs $v,w\in\mathbb{R}^n$. Then $k\geq n$.
\end{lemma}
\begin{proof}
  Suppose $k<n$ and obtain a contradiction. Let $S\subset\mathbb{R}^n$ be
  the unit sphere and consider $f:S\to\mathbb{R}$, $f(v)=\langle B(v,v),B(v,v)\rangle$.
  Let $v\in S$ be a minimum. Taking $v(t)=v\cos t+w\sin t$,
  
  $$
  0=df_v(w)=4\langle B(w,v),B(v,v)\rangle,\qquad(8)
  $$
  
  and, since $v''(0)=-v$,
  
  $$
  0\leq d^2f_v(w,w)=8\|B(w,v)\|^2-4\langle B(v,v),B(v,v)\rangle+4\langle B(w,w),B(v,v)\rangle.\qquad(9)
  $$
  
  Consider $L:T_vS\to\mathbb{R}^k$, $L(w)=B(v,w)$. Equation (8) yields
  $\langle L(w),B(v,v)\rangle=0$, hence $\dim L(T_vS)\leq k-1$. Because $k<n$, the
  kernel of $L$ has dimension $\geq 1$, so there exists $w_o\neq 0$, $w_o\perp v$,
  with $B(v,w_o)=0$. Introducing $w_o$ in (9),
  
  $$
  0\leq\langle B(w_o,w_o),B(v,v)\rangle-\|B(v,v)\|^2<\langle B(w_o,w_o),B(v,v)\rangle-(\sqrt{-b})^2\leq -b-(\sqrt{-b})^2=0,
  $$
  
  which produces the contradiction. $\blacksquare$
\end{proof}

\begin{corollary}[C. Tompkins]
  \label{cor:dc-ch10-3-4}
  \dcref{ch10:3.4}
  \uses{rem:dc-ch10-2-7}
  Suppose that $M$ is compact with zero curvature. If $k<n$, there does not exist
  an isometric immersion $f:M^n\to\mathbb{R}^{n+k}$; in particular, the flat torus
  $T^n$ (see \cref{ex:dc-ch1-2-7}, Chap. 1) cannot be immersed isometrically in $\mathbb{R}^{2n-1}$.
\end{corollary}

\begin{corollary}[B. O'Neill]
  \label{cor:dc-ch10-3-5}
  \dcref{ch10:3.5}
  Suppose that $\overline{M}$ is complete, simply connected and $K_{\overline{M}}\leq 0$.
  Let $M$ be compact, $\dim\overline{M}<2\dim M$ and $K_M\leq K_{\overline{M}}$.
  Then there does not exist an isometric immersion $f:M\to\overline{M}$.
\end{corollary}

\begin{remark}[Remark 3.6]
  \label{rem:dc-ch10-3-6}
  \dcref{ch10:3.6}
  \uses{thm:dc-ch10-3-1}
  The hypothesis that $M$ be compact is essential in \cref{thm:dc-ch10-3-1}, as is shown by
  the example of complete surfaces with $K\leq 0$ in $\mathbb{R}^3$. On the other
  hand, it is known that the hyperbolic plane $H^2$ (a complete surface with
  $K\equiv -1$) cannot be isometrically immersed in $\mathbb{R}^3$ (Theorem of
  Hilbert, see M. do Carmo [dC 2], p. 446). It is an open problem to know if the
  hyperbolic space $H^n$ can be immersed isometrically in $\mathbb{R}^{2n-1}$.
\end{remark}

\begin{remark}[Remark 3.7]
  \label{rem:dc-ch10-3-7}
  \dcref{ch10:3.7}
  \uses{thm:dc-ch10-3-1}
  Essentially the same proof used in \cref{thm:dc-ch10-3-1} can serve as a proof of the
  following theorem of S.S. Chern and N. Kuiper: Suppose that $M^n$ is compact and
  that, for each point $p\in M$, there exists a subspace of dimension $m$,
  $V^m\subset T_pM$, such that the sectional curvature with respect to the planes
  contained in $V$ is non-positive. If $k<m$, there does not exist an isometric
  immersion $f:M^n\to\mathbb{R}^{n+k}$. The historical importance of the article of
  Chern-Kuiper is that they stressed the fundamental fact that the existence of an
  isometric immersion $f:M^n\to\overline{M}^{n+k}$ implies a relation between the
  nullity of the second fundamental form of $f$ and the nullity of an intrinsic
  operator on $M$ defined from the curvature of $M$.
\end{remark}

\begin{remark}[Remark 3.8]
  \label{rem:dc-ch10-3-8}
  \dcref{ch10:3.8}
  \uses{thm:dc-ch10-3-1, rem:dc-ch10-2-7}
  The torus $T^{n+1}$ with the flat metric (Cf. \cref{ex:dc-ch1-2-7}, Chap. 1) contains the
  torus $T^n$ as a totally geodesic submanifold. This shows that the hypothesis
  that $\overline{M}$ be simply connected in \cref{thm:dc-ch10-3-1} is necessary.
\end{remark}

\section{Focal points and an extension of Rauch's Theorem}

The notion of a conjugate point to a given point $p\in M$ extends to the idea of a
focal point to a submanifold $N\subset M$. The idea is to consider variations
$f:(-\varepsilon,\varepsilon)\times[0,\ell]\to M$ of a geodesic
$\gamma:[0,\ell]\to M$ with $\gamma(0)=p\in N$ and $\gamma'(0)\in(T_pN)^\perp$,
satisfying:

1. The curve $t\mapsto f_s(t)$, $t\in[0,\ell]$, is a geodesic.
2. For all $s\in(-\varepsilon,\varepsilon)$, $f_s(0)=\alpha(s)\in N$ and
$A(s)=\frac{\partial f}{\partial t}(s,0)\in(T_{\alpha(s)}N)^\perp$.

It is clear that $J(t)=\frac{\partial f}{\partial s}(0,t)$ is a Jacobi field along
$\gamma$.

\begin{lemma}[Lemma 4.1]
  \label{lem:dc-ch10-4-1}
  \dcref{ch10:4.1}
  The Jacobi field $J$ constructed above satisfies the following properties:
  
  - i) $J(0)\in T_pN$,
  - ii) $J'(0)+S_{\gamma'(0)}(J(0))\in(T_pN)^\perp$,
  
  where $S_{\gamma'(0)}$ is the linear operator on $T_pN$ given by the second
  fundamental form of $N\subset M$. Conversely, if $J$ is a Jacobi field along the
  geodesic $\gamma$ with $\gamma(0)=p\in N$, $\gamma'(0)\in(T_pN)^\perp$, satisfying
  (i) and (ii), then there exists a variation $f$ of $\gamma$ satisfying (1) and (2),
  whose variational field is $J$.
\end{lemma}
\begin{proof}
  Clearly $J(0)=\frac{\partial f}{\partial s}(0,0)=\frac{d}{ds}\alpha(s)|_{s=0}\in T_pN$,
  which verifies (i). For (ii), let $v\in T_pN$; we show
  $\langle J'(0)+S_{\gamma'(0)}(J(0)),v\rangle=0$. Indeed,
  
  $$
  \Big\langle\frac{DJ}{dt}(0),v\Big\rangle=\Big\langle\frac{D}{ds}\frac{\partial f}{\partial t}(0,0),v\Big\rangle=\Big\langle\frac{DA}{ds}(0),v\Big\rangle.
  $$
  
  Since $A(s)$ is a field along $\alpha(s)$ with $\alpha'(0)=J(0)$,
  $\frac{DA}{ds}(0)=\overline\nabla_{J(0)}A(s)|_{s=0}$, so
  $\langle\frac{DA}{ds}(0),v\rangle=\langle-S_{A(0)}(J(0)),v\rangle=\langle-S_{\gamma'(0)}(J(0)),v\rangle$,
  which implies (ii).
  
  To prove the converse, let $s\mapsto\alpha(s)$ be a curve in $N$ with
  $\alpha(0)=p$, $\alpha'(0)=J(0)$. Choose a field $W$ along $\alpha$ with
  $W(0)=\gamma'(0)$ and $\frac{DW}{ds}(0)=\frac{DJ}{dt}(0)$, and write
  $W(s)=V(s)+U(s)$ with $V(s)\in(T_{\alpha(s)}N)^\perp$, $U(s)\in T_{\alpha(s)}N$.
  Define $f(s,t)=\exp_{\alpha(s)}(tV(s))$. One verifies $f$ satisfies (1) and (2),
  and that $\frac{DU}{ds}(0)=0$ (using that both $J$ and $\frac{\partial f}{\partial s}(0,t)$
  satisfy property (ii), and $U(0)=0$, $U(s)\in T_{\alpha(s)}N$), so that
  $\frac{\partial f}{\partial s}(0,t)=J(t)$. $\blacksquare$
\end{proof}

\begin{definition}[focal point]
  \label{def:dc-ch10-4-2}
  \dcref{ch10:4.2}
  Let $N\subset M$ be a submanifold of a Riemannian manifold $M$. The point $q\in M$
  is called a \emph{focal point} of $N$ if there exists a geodesic
  $\gamma:[0,\ell]\to M$, with $\gamma(0)=p\in N$, $\gamma'(0)\in(T_pN)^\perp$,
  $\gamma(\ell)=q$, and a non-zero Jacobi field $J$ along $\gamma$, satisfying (i),
  (ii) and with $J(\ell)=0$.
\end{definition}

\begin{example}[Example 4.3]
  \label{ex:dc-ch10-4-3}
  \dcref{ch10:4.3}
  If $S^{n-1}\subset S^n$ is the equator of $S^n$, that is,
  $S^{n-1}=\{x\in S^n;x=(x_1,\dots,x_n,0)\}$, then the north pole $(0,0,\dots,0,1)$
  and the south pole $(0,\dots,0,-1)$ are focal points of $S^{n-1}$ in $S^n$.
  
  To obtain a characterization of the focal points in terms of the exponential map,
  introduce the following notation. Let $T(M)$ be the tangent bundle of $M$,
  $\pi:T(M)\to M$ its projection. Similarly, $T(N)\to N$ denotes the tangent bundle
  of $N$, and $T(N)^\perp\to N$ denotes the normal bundle of the immersion
  $N\subset M$ (a point of $T(N)^\perp$ is a pair $(p,n)$, where $p\in N$,
  $n\in(T_pN)^\perp$). The exponential mapping can be thought of as a map
  $\exp:T(M)\to M$, defined by $\exp(p,v)=\exp_p(v)$. Since $T(N)^\perp\to N$ is
  contained in the restriction of $T(M)\to M$ to $N\subset M$, the exponential map
  can be restricted to $T(N)^\perp$; this restriction is denoted
  $\exp^\perp:T(N)^\perp\to M$. Observe that $\dim T(N)^\perp=\dim M$.
\end{example}

\begin{proposition}[Proposition 4.4]
  \label{prop:dc-ch10-4-4}
  \dcref{ch10:4.4}
  \uses{lem:dc-ch10-4-1}
  The point $q\in M$ is a focal point of $N\subset M$ if and only if it is a
  critical value of $\exp^\perp$.
\end{proposition}
\begin{proof}
  Suppose $q$ is a focal point of $N$. From \cref{lem:dc-ch10-4-1}, there exists a
  geodesic $\gamma:[0,\ell]\to M$ with $\gamma(0)=p\in N$, $\gamma(\ell)=q$,
  $\gamma'(0)\in(T_pN)^\perp$ and a variation $f$ of $\gamma$ satisfying (1) and (2)
  with $\frac{\partial f}{\partial s}(0,\ell)=0$. Then
  $w(s)=(\alpha(s),\ell A(s))$, $A(s)=\frac{\partial f}{\partial t}(s,0)$,
  $\alpha(s)=f_s(0)$, is a curve in $T(N)^\perp$ such that
  $\exp^\perp(w(s))=\exp_{\alpha(s)}\ell A(s)=f(s,\ell)$, $\exp^\perp(w(0))=q$, and
  $(d\exp^\perp)_{w(0)}(w'(0))=\frac{\partial f}{\partial s}(0,\ell)=0$. Then $w(0)$
  is a critical point of $\exp^\perp$, with $\exp^\perp(w(0))=q$.
  
  Conversely, suppose $q$ is a critical value of $\exp^\perp$. Then there exist
  $w_o$, $w_o'$ with $\exp^\perp(w_o)=q$ and $(d\exp^\perp)_{w_o}(w_o')=0$. Let
  $w(s)=(\sigma(s),\ell V(s))$ be a curve in $T(N)^\perp$ with $w(0)=w_o$,
  $w'(0)=w_o'$. Since $q=\exp_{\sigma(0)}(\ell V(0))$, there exists a geodesic
  $\gamma$ with $\gamma(0)=\sigma(0)$, $\gamma(\ell)=q$, $\gamma'(0)=V(0)$. Consider
  the variation $f(s,t)=\exp_{\sigma(s)}tV(s)$; it satisfies (1) and (2). From the
  lemma, $J(t)=\frac{\partial f}{\partial s}(0,t)$ is a Jacobi field satisfying (i)
  and (ii), with $J(\ell)=\frac{d}{ds}\exp^\perp w(s)|_{s=0}=(d\exp^\perp)_{w_o}(w_o')=0$,
  and therefore $q$ is a focal point of $N$. $\blacksquare$
\end{proof}

\begin{definition}[focal set]
  \label{def:dc-ch10-4-5}
  \dcref{ch10:4.5}
  Let $N\subset M$ be a submanifold of a Riemannian manifold $M$. The set of focal
  points of $N$ will be called the \emph{focal set} $F(N)\subset M$ of $N$.
\end{definition}

\begin{example}[Example 4.6]
  \label{ex:dc-ch10-4-6}
  \dcref{ch10:4.6}
  \uses{prop:dc-ch10-4-4}
  Let $N^2\subset\mathbb{R}^3$ be a regular surface in $\mathbb{R}^3$. The connected
  components of the focal set $F(N)\subset\mathbb{R}^3$ are called the "focal
  surfaces" (or surfaces of centers) in classical differential geometry. In general
  such subsets are not regular surfaces and can degenerate to points or curves. For
  example, if $N^2$ is a sphere $S^2$, then $F(S^2)$ is the center of $S^2$; if
  $N^2$ is the right circular cylinder $C$, $F(C)$ is the axis of $C$. A way of
  constructing geometrically the focal points of $N^2\subset\mathbb{R}^3$ is the
  following. Let $\mathbf{x}:U\subset\mathbb{R}^2\to N\subset\mathbb{R}^3$ be a
  parametrization of $N$ at $p\in N$, with coordinates $(u_1,u_2)$ and unit normal
  field $n=n(u_1,u_2)$; write $\exp^\perp((u_1,u_2),t)=\mathbf{x}(u_1,u_2)+tn$. By
  \cref{prop:dc-ch10-4-4}, $\mathbf{x}+tn$ is a focal point of $N$ if and only if the matrix
  of $d\exp^\perp$ is singular, which reduces to the matrix $(g_{ij}-th_{ij})$ being
  singular. Choosing $\mathbf{x}$ so that $(g_{ij})$ is the identity at $p$,
  $\mathbf{x}+tn$ is a focal point if and only if $\frac{1}{t}$ is an eigenvalue of
  $(h_{ij})$, that is, $\frac{1}{t}$ is one of the principal curvatures of $N^2$ at
  $p$. In summary, to construct a connected component of $F(N)$, fix a principal
  curvature $k_i$ and mark on the normal at each $p\in N$, in the direction $n(p)$,
  a length equal to $1/k_i(p)$ (the center of the osculating circle in the normal
  section of $N$ in the principal direction $e_i$, which justifies the name "surface
  of centers").
  
  We now pass to an extension of Rauch's Theorem in which the notion of focal point
  replaces that of conjugate point.
\end{example}

\begin{definition}[focal point free]
  \label{def:dc-ch10-4-7}
  \dcref{ch10:4.7}
  Let $\gamma:[0,a]\to M$ be a geodesic in $M$. Let $B_\varepsilon(0)\subset\gamma'(0)^\perp$
  be a ball of radius $\varepsilon$ and center at the origin $0$, contained in the
  orthogonal complement of $\gamma'(0)$. We say that $\gamma$ is \emph{focal point free}
  at $(0,a]$ if there exists some $\varepsilon>0$ such that $\gamma$ has no focal
  points relative to the submanifold $\Sigma_\varepsilon=\exp_{\gamma(0)}(B_\varepsilon(0))$.
  
  Observe that since $\Sigma_\varepsilon$ is geodesic at $\gamma(0)$, any Jacobi field $J$
  along $\gamma$, with $J(0)\neq 0$ and $J'(0)=0$, automatically satisfies
  $S_{\gamma'(0)}(J(0))=0$.
\end{definition}

\begin{lemma}[Index Lemma for focal points]
  \label{lem:dc-ch10-4-8}
  \dcref{ch10:4.8}
  Let $\gamma:[0,a]\to M^n$ be a geodesic which is focal point free on $(0,a]$. Let
  $J$ be a Jacobi field along $\gamma$, with $\langle J,\gamma'\rangle=0$, and let
  $V$ be a piecewise differentiable vector field along $\gamma$. Suppose that
  $J'(0)=0$ and $J(t_o)=V(t_o)$, $t_o\in(0,a]$. Then
  
  $$
  I_{t_o}(J,J)\leq I_{t_o}(V,V),
  $$
  
  and the equality occurs if and only if $V=J$ on $[0,t_o]$.
\end{lemma}
\begin{proof}
  Let $\{J_1,\dots,J_{n-1}\}$ be a basis of the vector space of Jacobi
  fields $J$ such that $J'(0)=0$, $\langle J,\gamma'\rangle=0$. The fact that
  $\gamma$ is focal point free on $(0,a]$ implies that, for each $t\in(0,a]$,
  $\{J_1(t),\dots,J_{n-1}(t)\}$ is a basis for $(\gamma'(t))^\perp$. Starting from
  there, the proof follows in a manner entirely analogous to the Index Lemma. $\blacksquare$
\end{proof}

\begin{theorem}[extension of Rauch to focal points]
  \label{thm:dc-ch10-4-9}
  \dcref{ch10:4.9}
  \uses{lem:dc-ch10-4-8, thm:dc-ch10-2-3}
  Let $\gamma:[0,a]\to M^n$ and $\tilde\gamma:[0,a]\to\widetilde{M}^{n+k}$ be
  geodesics, with the same velocity, and let $J$ and $\tilde J$ be Jacobi fields
  along $\gamma$ and $\tilde\gamma$, such that
  
  $$
  0=J'(0)=\tilde J'(0),\quad\langle J(0),\gamma'(0)\rangle=\langle \tilde J(0),\tilde\gamma'(0)\rangle,\quad|J(0)|=|\tilde J(0)|.
  $$
  
  Assume that $\tilde\gamma$ is focal point free on $(0,a]$ and that, for all $t$
  and all $x\in T_{\gamma(t)}(M)$, $\tilde x\in T_{\tilde\gamma(t)}(\widetilde{M})$,
  we have
  
  $$
  \widetilde{K}(\tilde x,\tilde\gamma'(t))\geq K(x,\gamma'(t)).
  $$
  
  Then $|\tilde J|\leq|J|$. In addition, if for some $t_o\in(0,a]$ we have
  $|\tilde J(t_o)|=|J(t_o)|$, then $\widetilde{K}(\tilde J(t),\tilde\gamma'(t))=K(J(t),\gamma'(t))$
  for all $t\in[0,t_o]$.
\end{theorem}
\begin{proof}
  The proof is analogous to that of \cref{thm:dc-ch10-2-3} with the following
  modifications: $\frac{|J|^2}{|\tilde J|^2}$ is well-defined, since $\tilde\gamma$
  is focal point free, and $I_{t_o}(\phi(U),\phi(U))\geq I_{t_o}(\tilde U,\tilde U)$
  by \cref{lem:dc-ch10-4-8}. $\blacksquare$
  
  Other useful extensions of the Rauch comparison theorem can be found in F. Warner,
  "Extensions of the Rauch comparison theorem to submanifolds", Trans. A.M.S. 122
  (1966), 341-356, and in E. Heintze and H. Karcher, "A general comparison theorem
  with applications to volume estimates for submanifolds", Ann. Sci. Ecole Norm.
  Sup. 11 (1978), 451-470.
  
  Rauch's theorem admits an extremely important global generalization, called the
  Theorem of Toponogov. One of its versions: let $M$ be a Riemannian manifold which
  is complete with sectional curvature $K\geq H$. Let $\gamma_1$ and $\gamma_2$ be
  normalized geodesic segments in $M$ with $\gamma_1(0)=\gamma_2(0)$. Denote by
  $M^2(H)$ a manifold of dimension two with constant curvature $H$. Assume that
  $\gamma_1$ is minimizing and that, if $H>0$, $\ell(\gamma_2)\leq\pi/\sqrt{H}$.
  Consider on $M^2(H)$ two normalized geodesics $\bar\gamma_1$, $\bar\gamma_2$ with
  $\bar\gamma_1(0)=\bar\gamma_2(0)$, $\ell(\gamma_i)=\ell(\bar\gamma_i)=\ell_i$,
  $i=1,2$, and $\angle(\bar\gamma_1'(0),\bar\gamma_2'(0))=\angle(\gamma_1'(0),\gamma_2'(0))$.
  Then
  
  $$
  d(\gamma_1(\ell_1),\gamma_2(\ell_2))\leq d(\bar\gamma_1(\ell_1),\bar\gamma_2(\ell_2)).
  $$
  
  A proof can be found in Cheeger and Ebin [CE]. The Theorem of Toponogov is an
  essential tool for the study of the relation between topology and curvature; one
  of the culmination points, the sphere theorem, is presented in Chapter 13.
\end{proof}
